\section{Resultados}
\subsection{Avaliações Individuais}
\begin{table}[H]
\centering
\begin{tabular}{ | c | c |}
  \hline
  função mais comum nos padrões & free \\
  \hline
  maior suporte & 16,43\% \\
  \hline
  suporte médio & 1,98\% \\
  \hline
  suporte mediano & 1,22\% \\
  \hline
  menor suporte & 0.70\% \\
  \hline
  desvio padrão suporte & 2,16\% \\
  \hline
  padrão de uso com maior suporte & malloc, free\\
  \hline
  padrão de uso com menor suporte & memset, calloc, \\
   & snprintf,strcmp\\
  \hline
  maior tamanho de padrão de uso & 4 \\
  \hline
  tamanho médio dos padrões de uso & 2,6 \\
  \hline
  tamanho mediano dos padrões de uso & 2\\
  \hline
  menor tamanho de padrão de uso & 2\\
  \hline
  desvio padrão do tamanho de padrão de uso & 0,7 \\
  \hline
  maior padrão de uso detectado & calloc, memcpy, \\
  &  free, malloc\\
  \hline
  menor padrão de uso detectado & strstr, strlen\\
  \hline
  número de padrões detectados & 150 \\
  \hline
\end{tabular}
\caption{Análise estatística dos padrões detectados via mineração de regras de associação}
\label{desc_item}
\end{table}

\begin{table}[H]
\centering
  \begin{tabular}{ | c | c |}
\hline
função mais comum nos padrões detectados & malloc \\
\hline
maior suporte & 16,43\% \\
\hline
suporte médio & 2,13\% \\
\hline
suporte mediano & 1,02\% \\
\hline
menor suporte & 0,00\% \\
\hline
desvio padrão do suporte & 4,70\% \\
\hline
padrão de uso com maior suporte & free, malloc \\
\hline
padrão de uso com menor suporte & sprintf, strcpy \\
&  memcpy, free, strcmp \\
\hline
maior tamanho de padrão de uso & 8\\
\hline
tamanho médio dos padrões de uso & 2.65 \\
\hline
tamanho mediano dos padrões de uso & 2 \\
\hline
menor tamanho de padrão de uso & 2\\
\hline
desvio padrão do tamanho de padrão de uso & 1.03 \\
\hline
maior padrão de uso detectado & fread, strlen, memcpy\\
& memset, free, strchr\\
& fgetc, strcmp\\
\hline
menor padrão de uso detectado & inb\_p, outb\_p\\
\hline
número de padrões detectados & 82\\
\hline
\end{tabular}
\caption{Análise estatística dos padrões detectados via clustering}
\label{desc_clustering}
\end{table}

\begin{table}[H]
\centering
\begin{tabular}{ | c | c |}
  \hline
  função mais comum nos padrões detectados & close \\
  \hline
  maior suporte & 4,70\% \\
  \hline
  suporte médio &  0,21\%\\
  \hline
  suporte mediano & 0,11\% \\
  \hline
  menor suporte & 0,05\% \\
  \hline
  desvio padrão do suporte & 0,4\% \\
  \hline
  padrão de uso com maior suporte & pthread\_mutex\_unlock, \\
   & pthread\_mutex\_lock \\
  \hline
  padrão de uso com menor suporte & fork, \_exit \\
  \hline
  maior tamanho de padrão de uso & 5\\
  \hline
  tamanho médio dos padrões de uso & 2,45 \\
  \hline
  tamanho mediano dos padrões de uso & 2 \\
  \hline
  menor tamanho de padrão de uso & 2\\
  \hline
  desvio padrão do tamanho de padrão de uso & 0,73 \\
  \hline
  maior padrão de uso detectado & strcmp, strstr\\
   & strdup, malloc, calloc\\
  \hline
  menor padrão de uso detectado & vprintf, putchar\\
  \hline
  número de padrões detectados & 111 \\
  \hline
\end{tabular}
\caption{Análise estatística dos padrões detectados via modelo probabilístico}
\label{desc_pam}
\end{table}

\subsection{Avaliações Individuais}

\begin{table}[H]
\centering
\begin{tabular}{ | c | c | c | }
  \hline
  Combinação & Indice de Jacquard \\
  \hline
  Regras de Associação - Modelo Probabilístico & 1,14\% \\
  \hline
  Regras de Associação - Clustering &  18,31\% \\
  \hline
  Modelo Probabilístico - Clustering & 2,52\% \\
  \hline
  Todos & 0,7 \% \\
  \hline
\end{tabular}
\caption{ Comparações entre os padrões dos métodos com \% da interseção dos
  dados entre os conjuntos }
\label{comp}
\end{table}
